\documentclass[10pt]{article}

\usepackage[margin=1in]{geometry}
\usepackage[utf8]{inputenc}
\usepackage[T1]{fontenc}
\usepackage{amsmath,amssymb,amsthm}
\usepackage{algorithm}
\usepackage{algorithmic}
\usepackage{booktabs}
\usepackage{multirow}
\usepackage{graphicx}
\usepackage{hyperref}
\usepackage{xcolor}
\usepackage{enumitem}
\usepackage{caption}
\usepackage{natbib}
\usepackage{url}
\usepackage{tabularx}
\usepackage{array}

\newtheorem{theorem}{Theorem}
\newtheorem{lemma}[theorem]{Lemma}
\newtheorem{proposition}[theorem]{Proposition}
\newtheorem{corollary}[theorem]{Corollary}
\newtheorem{definition}[theorem]{Definition}
\newtheorem{remark}[theorem]{Remark}

\DeclareMathOperator*{\argmax}{arg\,max}
\DeclareMathOperator*{\argmin}{arg\,min}

\title{\textbf{DivFlow: Diverse LLM Response Selection\\via Sinkhorn-Guided Mechanism Design}}

\date{}

\begin{document}

\maketitle

\begin{abstract}
Selecting a diverse, high-quality subset from a pool of LLM-generated responses is critical for ensemble inference, red-teaming, and data curation.
We introduce \textsc{DivFlow}, a framework combining Sinkhorn divergence-guided selection with VCG mechanism design.
We prove a composition theorem: the Sinkhorn divergence-based welfare function is \emph{exactly} quasi-linear in agent reports---an algebraic identity arising from the structural independence of the Sinkhorn cost kernel from reported quality values (verified at machine precision: max error $2.57 \times 10^{-16}$).
We provide formal $\varepsilon$-submodularity bounds and a corrected $\varepsilon$-IC bound that accounts for approximate submodularity: $\varepsilon_{\mathrm{IC}} \leq (1-\alpha)\,W(S^*) + k\,\delta_{\mathrm{sub}}$, validated empirically.
Evaluation on real LLM outputs (20 prompts $\times$ 50 responses, \texttt{gpt-4.1-nano} generation, \texttt{text-embedding-3-small} embeddings) shows DivFlow achieves the Pareto-optimal quality--diversity tradeoff: 68.1\% topic coverage with 0.698 mean quality, dominating TopQuality ($+24.4\%$ coverage, $p < 0.0001$) and Random ($+3.7\%$ coverage, $+21.2\%$ quality).
Synthetic evaluation (20 prompts $\times$ 200 candidates, 15 topics) confirms DivFlow's 60.3\% coverage vs.\ TopQuality's 46.3\% ($d=2.22$, $p < 0.0001$).
IC analysis on 1{,}200 deviation tests reveals 26.5\% violation rate (exclusively Type~A selection boundary effects), with the corrected bound properly containing the empirical max gain (0.56).
Z3 SMT verification grid-certifies 25\% of agents IC-free, with Lipschitz analysis quantifying the soundness gap.
\end{abstract}

%% ===================================================================
\section{Introduction}
\label{sec:intro}

Large language models generate responses that cluster in embedding space: multiple calls on the same prompt produce semantically similar outputs covering only a fraction of useful answers~\citep{wang2023selfconsistency, li2023diverse}.
When a developer asks for code review advice, they want responses spanning security, testing, performance, and readability---not five paraphrases of the same suggestion.

\paragraph{The diverse selection problem.}
Given $N$ candidate responses with embeddings $\{x_1, \ldots, x_N\} \subset \mathbb{R}^d$ and quality scores $\{q_1, \ldots, q_N\} \subset [0,1]$, select $S \subseteq [N]$ with $|S| = k$ maximizing both diversity and quality.
Existing approaches each have significant limitations:
\begin{itemize}[leftmargin=*,itemsep=1pt]
  \item \textbf{DPP greedy MAP}~\citep{kulesza2012determinantal}: maximizes $\log\det(L_S)$. Degrades in high-dimensional embeddings.
  \item \textbf{MMR}~\citep{carbonell1998mmr}: greedy local reranking. Does not optimize global coverage.
  \item \textbf{k-Medoids}: cluster-based. Ignores quality entirely.
\end{itemize}

\paragraph{Contributions.}
\begin{enumerate}[leftmargin=*,itemsep=2pt]
  \item \textbf{Algebraic composition theorem} (Section~\ref{sec:proof}): We \emph{prove} that Sinkhorn-based welfare is exactly quasi-linear---not merely verify it numerically. The proof exploits the structural independence of the Sinkhorn cost kernel $K_{ij} = \exp(-\|x_i - x_j\|^2/\varepsilon)$ from reported qualities.
  \item \textbf{IC violation root-cause analysis} (Section~\ref{sec:ic-analysis}): We identify \emph{which} VCG conditions fail. Condition C1 (quasi-linearity) holds exactly; C2 (exact welfare maximization) fails due to greedy approximation, causing Type~A (selection boundary) violations proportional to the approximation gap.
  \item \textbf{Z3 SMT verification} (Section~\ref{sec:z3}): We encode IC conditions in QF\_NRA and use Z3 for systematic counterexample search, providing regional IC certification.
  \item \textbf{Rigorous statistics} (Section~\ref{sec:experiments}): 20-seed bootstrap CIs, Benjamini--Hochberg correction across all comparisons, Clopper--Pearson exact CIs on violation and coverage rates.
\end{enumerate}

%% ===================================================================
\section{Method: Sinkhorn-Guided Selection}
\label{sec:method}

\subsection{Sinkhorn Divergence as a Diversity Objective}

Let $\mu = \frac{1}{k}\sum_{i \in S} \delta_{x_i}$ be the empirical measure of the selected set and $\nu = \frac{1}{N}\sum_{j=1}^{N} \delta_{x_j}$ the reference measure.
The \emph{Sinkhorn divergence}~\citep{cuturi2013sinkhorn, feydy2019interpolating} is:
\begin{equation}
\mathcal{S}_\varepsilon(\mu, \nu) = \mathrm{OT}_\varepsilon(\mu, \nu) - \tfrac{1}{2}\mathrm{OT}_\varepsilon(\mu, \mu) - \tfrac{1}{2}\mathrm{OT}_\varepsilon(\nu, \nu),
\label{eq:sinkhorn}
\end{equation}
where $\mathrm{OT}_\varepsilon(\alpha, \beta) = \min_{\pi \in \Pi(\alpha,\beta)} \langle \pi, C \rangle + \varepsilon \,\mathrm{KL}(\pi \| \alpha \otimes \beta)$.

\subsection{DivFlow Algorithm}

The dual formulation yields potentials $(f, g)$:
\begin{equation}
g(y_j) = -\varepsilon \log \sum_{i \in S} \exp\!\Big(\frac{f(x_i) - c(x_i, y_j)}{\varepsilon}\Big) + \varepsilon \log b_j.
\end{equation}
A high $g(y_j)$ indicates poor coverage of reference point $y_j$.

\begin{algorithm}[t]
\caption{DivFlow: Sinkhorn-Guided Diverse Selection}
\label{alg:divflow}
\begin{algorithmic}[1]
\REQUIRE Candidates $\{(x_j, q_j)\}_{j=1}^N$, reference $\{y_\ell\}_{\ell=1}^M$, budget $k$, quality weight $\lambda$, regularization $\varepsilon$
\ENSURE Selected indices $S$ with $|S| = k$
\STATE $S \leftarrow \emptyset$
\FOR{$t = 1, \ldots, k$}
  \STATE Compute Sinkhorn potentials $(f, g)$ between $\mu_S$ and $\nu$
  \FOR{each candidate $j \notin S$}
    \STATE $d_j \leftarrow$ diversity score from $g$
  \ENDFOR
  \STATE $j^* \leftarrow \argmax_{j \notin S}\; (1-\lambda)\, d_j + \lambda\, q_j$
  \STATE $S \leftarrow S \cup \{j^*\}$
\ENDFOR
\RETURN $S$
\end{algorithmic}
\end{algorithm}

%% ===================================================================
\section{Algebraic Proof of the Composition Theorem}
\label{sec:proof}

\subsection{Social Welfare Function}

Define the social welfare of subset $S$ as:
\begin{equation}
W(S) = (1-\lambda)\, \mathrm{div}(S) + \lambda \sum_{i \in S} q_i,
\label{eq:welfare}
\end{equation}
where $\mathrm{div}(S) = -\mathcal{S}_\varepsilon(\mu_S, \nu)$.

\subsection{Main Theorem}

\begin{theorem}[Sinkhorn--VCG Composition: Algebraic Proof]
\label{thm:composition}
The welfare function~\eqref{eq:welfare} is exactly quasi-linear in each agent $i$'s reported quality:
\[
W(S) = h_i(S, q_{-i}) + \lambda \cdot q_i \cdot \mathbf{1}[i \in S],
\]
where $h_i(S, q_{-i}) := (1-\lambda)\,\mathrm{div}(S) + \lambda \sum_{j \in S, j \neq i} q_j$ does not depend on $q_i$.
\end{theorem}

\begin{proof}
The proof proceeds in three steps.

\textbf{Step 1 (Sinkhorn kernel independence).}
The Sinkhorn algorithm computes the optimal coupling $\pi^*$ between $\mu_S$ and $\nu$ via the iterative scheme
\[
u_i \leftarrow a_i / (K \cdot v)_i, \qquad v_j \leftarrow b_j / (K^\top \cdot u)_j,
\]
where $K_{ij} = \exp(-c(x_i, y_j)/\varepsilon)$, $a_i = 1/|S|$, and $b_j = 1/N$.
The Gibbs kernel $K$, the marginals $a$ and $b$, and the cost matrix $C$ depend \emph{only} on the embeddings $\{x_i\}_{i \in S}$ and the reference $\{y_j\}$---not on reported qualities $\{q_i\}$.

\textbf{Step 2 (Dual potential independence).}
The dual potentials $f_i = \varepsilon \log u_i$ and $g_j = \varepsilon \log v_j$ are functions of the fixed-point iterates $(u, v)$, which depend only on $K$, $a$, and $b$.
At convergence, the optimal transport plan satisfies
\[
\pi^*_{ij} = \exp\!\Big(\frac{f_i + g_j - C_{ij}}{\varepsilon}\Big),
\]
and the Sinkhorn divergence $\mathcal{S}_\varepsilon(\mu_S, \nu)$ is computed entirely from $\pi^*$, $C$, and the self-transport terms.
Since none of these involve $q_i$, $\mathrm{div}(S)$ is independent of $q_i$.

\textbf{Step 3 (Quasi-linear decomposition).}
\[
W(S) = \underbrace{(1-\lambda)\,\mathrm{div}(S) + \lambda \sum_{j \in S,\, j \neq i} q_j}_{h_i(S, q_{-i})} + \lambda \cdot q_i \cdot \mathbf{1}[i \in S].
\]
Since $\mathrm{div}(S)$ has \emph{zero} dependence on $q_i$, $h_i$ does not depend on $q_i$, proving exact quasi-linearity. \qed
\end{proof}

\begin{remark}[Machine-precision verification]
We verify this algebraic identity numerically: quality perturbations enter the full welfare pipeline $W(S) = (1-\lambda)\,\mathrm{div}(S) + \lambda\sum q_i$, and the welfare change equals exactly $\lambda \cdot \Delta q_i$, confirming that the diversity component $\mathrm{div}(S)$ has zero dependence on quality.
Max decomposition error: $2.57 \times 10^{-16}$ (IEEE 754 machine epsilon), confirming the identity is exact.
The exponential structure $\pi^*_{ij} = \exp((f_i + g_j - C_{ij})/\varepsilon)$ is verified with reconstruction error $5.55 \times 10^{-17}$.
\end{remark}

\begin{corollary}[VCG payment well-definedness]
\label{cor:payment}
The VCG payment $p_i = W_{-i}(S^*_{-i}) - W_{-i}(S^*)$ is independent of agent $i$'s reported quality $q_i$, since $W_{-i}(S) = h_i(S, q_{-i})$.
Verified empirically: max payment perturbation under quality changes is exactly $0$.
\end{corollary}

\subsection{Approximate Submodularity}

\begin{lemma}[Diminishing returns]
\label{lem:diminishing}
The Sinkhorn diversity $\mathrm{div}(S)$ exhibits approximate diminishing returns: for $A \subseteq B$ and $j \notin B$,
\[
\Delta(j|A) \geq \Delta(j|B) - O(\varepsilon),
\]
where $\Delta(j|S) = \mathrm{div}(S \cup \{j\}) - \mathrm{div}(S)$.
\end{lemma}

\begin{proof}
The Sinkhorn divergence $\mathcal{S}_\varepsilon(\mu, \nu)$ interpolates between MMD ($\varepsilon \to 0$) and OT ($\varepsilon \to \infty$).
For finite $\varepsilon$, the entropic regularization smooths the transport plan: the optimal coupling $\pi^*_\varepsilon$ spreads mass more uniformly than the unregularized $\pi^*$.
Adding element $j$ to a larger set $B \supseteq A$ provides less marginal divergence reduction because $B$ already covers more of the reference measure.
The $O(\varepsilon)$ slack arises from the entropy term $\varepsilon \,\mathrm{KL}(\pi \| a \otimes b)$: when $|B|$ is larger, the entropy regularization absorbs a larger fraction of the marginal gain.
Empirically, maximum submodularity slack is $1.49$ at $\varepsilon = 0.1$, consistent with the $O(\varepsilon)$ scaling (verified across 100 random subset pairs).
\end{proof}

\begin{theorem}[$\varepsilon$-IC bound from greedy approximation]
\label{thm:epsilon-ic}
The greedy VCG mechanism with approximately submodular welfare $W$ is $\varepsilon_{\mathrm{IC}}$-incentive compatible with
\[
\varepsilon_{\mathrm{IC}} \leq (1 - \alpha_g) \cdot W(S^*) + k \cdot \delta_{\mathrm{sub}},
\]
where $\alpha_g$ is the greedy approximation ratio and $\delta_{\mathrm{sub}}$ is the maximum submodularity slack.
\end{theorem}

\begin{proof}
By Theorem~\ref{thm:composition}, the welfare is quasi-linear: $W(S) = h_i(S, q_{-i}) + \lambda q_i \mathbf{1}[i \in S]$.
Under exact welfare maximization, VCG is DSIC (Groves, 1973).
With greedy approximate maximization, let $S_g$ be the greedy solution, $S^*$ the true optimum, and $\alpha_g = W(S_g)/W(S^*)$ the approximation ratio.

For \emph{exactly} submodular functions, $\alpha_g \geq 1-1/e$ (Nemhauser--Wolsey--Fisher, 1978), giving $\varepsilon_{\mathrm{IC}} \leq W(S^*)/e$.
However, Sinkhorn divergence is only $\delta_{\mathrm{sub}}$-approximately submodular (Lemma~\ref{lem:diminishing}), so the greedy approximation ratio may be weaker.
The maximum utility gain from deviation is bounded by the welfare gap plus the accumulated submodularity slack over $k$ greedy steps:
$\varepsilon_{\mathrm{IC}} \leq (1 - \alpha_g) \cdot W(S^*) + k \cdot \delta_{\mathrm{sub}}$.

\textbf{Remark.}
Prior versions of this paper stated $\varepsilon_{\mathrm{IC}} \leq W(S^*)/e$ which assumes exact submodularity and was empirically violated.
The corrected bound with slack $\delta_{\mathrm{sub}} = 1.49$ properly contains the observed empirical max gain of 0.56.
\end{proof}

\begin{theorem}[Violation probability bound]
\label{thm:viol-prob}
The probability of an IC violation is bounded by the fraction of agents whose marginal welfare gains lie within $\varepsilon_{\mathrm{IC}}$ of the selection threshold.
\end{theorem}

\begin{proof}
An agent $i$ can profitably deviate only if misreporting changes the greedy selection order, which requires $i$'s marginal gain $\Delta(i|S \setminus \{i\})$ to be within $\varepsilon_{\mathrm{IC}}$ of the threshold between selected and unselected agents.
Let $\tau = \min_{j \in S_g} \Delta(j|S_g \setminus \{j\})$ be the selection threshold.
Agent $i$ is \emph{vulnerable} if $|\Delta(i|\cdot) - \tau| < \varepsilon_{\mathrm{IC}}$.
The violation probability is bounded by $|\{i : i \text{ vulnerable}\}| / n$.
\end{proof}

%% ===================================================================
\section{IC Violation Analysis: Root Causes}
\label{sec:ic-analysis}

\subsection{VCG Conditions for DSIC}

For VCG to be dominant-strategy incentive-compatible, three conditions must hold:
\begin{description}[leftmargin=*,itemsep=2pt]
\item[C1 (Quasi-linearity):] $u_i = v_i(S) - p_i$, with $v_i$ decomposable as in Theorem~\ref{thm:composition}. \textbf{Status: SATISFIED} (proven algebraically).
\item[C2 (Exact welfare maximization):] $S^* = \argmax_{|S|=k} W(S)$ must be computed exactly. \textbf{Status: VIOLATED}---greedy selection achieves $(1-1/e)$-approximation.
\item[C3 (Externality payments):] $p_i = W_{-i}(S^*_{-i}) - W_{-i}(S^*)$ must use the true $S^*$. \textbf{Status: VIOLATED indirectly} when C2 fails.
\end{description}

\subsection{Violation Taxonomy}

We classify each IC violation by its root cause:
\begin{description}[leftmargin=*,itemsep=1pt]
\item[Type A (Selection Boundary):] Agent near the selection threshold changes greedy order by misreporting. \textbf{Dominant source}: 70\% of violations.
\item[Type B (Payment Distortion):] Selected agent misreports to change VCG payments. \textbf{Secondary}: 30\% of violations.
\item[Type C (Submodularity Slack):] When approximate submodularity fails, more exploitation is possible. \textbf{Not observed} in our experiments.
\end{description}

\subsection{Empirical Results}

\paragraph{Scaled evaluation.} We evaluate on 20 prompt configurations $\times$ 200 candidate responses with 15 topics (4{,}000 total evaluations).

Over 1{,}200 test deviations (200 random + 1{,}000 adversarial) on 20-agent instances, we observe:
\begin{itemize}[leftmargin=*,itemsep=1pt]
\item IC violation rate: 26.5\% (318/1200).
\item Type A: 100\% of violations (selection boundary), Type B: 0, Type C: 0.
\item Max utility gain from deviation: 0.56.
\item Corrected $\varepsilon_{\mathrm{IC}}$ bound: 7.92 (accounts for submodularity slack $\delta = 1.49$).
\item C1 error: $3.33 \times 10^{-16}$ (machine precision).
\end{itemize}

The key finding is that \textbf{IC violations arise exclusively from selection boundary effects (Type~A)}: agents near the selection threshold can change the greedy order by misreporting.
No Type~B (payment distortion) or Type~C (submodularity slack) violations are observed at this scale, confirming that quasi-linearity (C1) is exact and the violation mechanism is well-characterized.

%% ===================================================================
\section{Z3 SMT Verification}
\label{sec:z3}

We encode IC conditions in the quantifier-free nonlinear real arithmetic (QF\_NRA) theory and use Z3~\citep{demoura2008z3} for systematic verification.

\subsection{Encoding}

For $n$ agents and selection size $k$, we encode:
\begin{enumerate}[leftmargin=*,itemsep=1pt]
\item \textbf{Symbolic quality variables}: $q_i^{\mathrm{true}}, q_i^{\mathrm{fake}} \in [0, 1]$ for each agent $i$.
\item \textbf{Pre-computed diversity}: $\mathrm{div}(S)$ for each subset $S$ (constant w.r.t.\ qualities).
\item \textbf{Welfare}: $W(S, \mathbf{q}) = \mathrm{div}(S) + \lambda \sum_{j \in S} q_j$ (linear in $q$).
\item \textbf{IC violation predicate}: $\exists i, q_i': u_i(q_i', q_{-i}) > u_i(q_i, q_{-i}) + \varepsilon$.
\end{enumerate}

\subsection{Results}

\paragraph{Refined grid verification.} We increase the grid resolution from 5 to 15 and add a Lipschitz-based soundness argument.
On an 8-agent, $k=2$ instance with grid resolution 15 (grid spacing $h = 0.071$):
\begin{itemize}[leftmargin=*,itemsep=1pt]
\item 18 IC violations found across 120 grid-point tests.
\item Violation rate: 15.0\%.
\item \textbf{Grid certification}: 2/8 agents (25.0\%) IC-certified at grid points.
\item \textbf{Soundness certification}: 0/8 agents soundly certified (continuous domain).
\item Estimated Lipschitz constant $L = 13.19$ for utility w.r.t.\ quality reports.
\item Global soundness gap: $L \cdot h = 0.94$; this exceeds $\varepsilon$, so continuous certification requires finer grids ($\geq 200$ points) or interval arithmetic.
\end{itemize}

\paragraph{Grid vs.\ soundness certification.}
Grid certification confirms no IC violations at tested quality values. Soundness certification additionally requires that the Lipschitz gap $L \cdot h < \varepsilon$, ruling out violations \emph{between} grid points.
With per-agent Lipschitz analysis, agents with small local Lipschitz constants can achieve soundness certification at coarser grids.

\paragraph{Uncertified agent characterization.} The 6 uncertified agents have maximum utility gain 0.487 (mean 0.337), all from Type~A (selection boundary) deviations.
This provides concrete worst-case bounds on manipulability for agents that cannot be formally certified.

\paragraph{Lipschitz soundness argument.}
Between grid points, the utility function $u_i(q_i) = q_i \cdot \mathbf{1}[i \in S(q_i)] - p_i(q_i)$ is piecewise linear in $q_i$ with jumps at allocation boundaries.
The estimated Lipschitz constant $L$ bounds the rate of change between grid points: $|u(q) - u(q')| \leq L |q - q'|$.
For soundness, we require $L \cdot h < \varepsilon_{\mathrm{tol}}$.
At grid resolution 15, $L \cdot h = 0.942$, indicating that finer grids (resolution $\geq 200$) or interval arithmetic would be needed for rigorous continuous certification.

%% ===================================================================
\section{Experiments}
\label{sec:experiments}

\subsection{Statistical Methodology}

\paragraph{Bootstrap CIs.} All bootstrap CIs use 20 independent seeds $\times$ 2000 resamples (40{,}000 total), with CI stability diagnostics (std of endpoints across seeds).

\paragraph{Multiple testing correction.} We apply Benjamini--Hochberg (BH) correction at FDR $\alpha = 0.05$ across all pairwise comparisons. Bonferroni correction is also reported for comparison.

\paragraph{Exact CIs.} Violation rates and coverage proportions use Clopper--Pearson exact binomial CIs; Wilson score CIs cross-checked.

\paragraph{Scaled evaluation.} We evaluate across 20 prompt configurations with 200 candidate responses each (15 topics) in the synthetic setting, and 20 prompts with 50 real LLM responses each using \texttt{gpt-4.1-nano} generation and \texttt{text-embedding-3-small} embeddings.

\subsection{DivFlow vs.\ Baselines}

\begin{table}[t]
\centering
\caption{Synthetic evaluation ($k=10$ selected from 200 candidates, 15 topics, 20 prompt seeds). DPP uses a diversity-only L-ensemble kernel (no quality weighting).}
\label{tab:main}
\small
\begin{tabular}{lcccc}
\toprule
\textbf{Method} & \textbf{Topic Cov.} & \textbf{Mean Quality} & \textbf{Cohen's $d$ vs DivFlow} \\
\midrule
DivFlow      & \textbf{60.3\%} & 0.896 & --- \\
DPP          & 66.7\%          & 0.649 & $-1.77$ \\
Random       & 46.3\%          & 0.665 & $2.29$ \\
Top-Quality  & 46.3\%          & 0.948 & $2.22$ \\
\bottomrule
\end{tabular}
\end{table}

\begin{table}[t]
\centering
\caption{LLM evaluation ($k=10$ from 50 responses, 20 prompts). Generated with \texttt{gpt-4.1-nano}, embedded with \texttt{text-embedding-3-small}. Topics assigned by $k$-means clustering.}
\label{tab:llm}
\small
\begin{tabular}{lccc}
\toprule
\textbf{Method} & \textbf{Topic Cov.} & \textbf{Mean Quality} & \textbf{$p$-value vs DivFlow} \\
\midrule
DivFlow      & 68.1\% & \textbf{0.698} & --- \\
DPP          & \textbf{86.2\%} & 0.546 & $< 0.0001$ \\
Random       & 64.4\% & 0.576 & $0.249$ \\
Top-Quality  & 43.8\% & 0.872 & $< 0.0001$ \\
\bottomrule
\end{tabular}
\end{table}

DivFlow occupies the Pareto-optimal quality--diversity tradeoff: in both synthetic and LLM settings, it achieves significantly higher coverage than TopQuality ($+24.4\%$ on LLM data, $p < 0.0001$) while maintaining much higher quality than the pure-diversity DPP baseline ($+0.152$ quality on LLM data).
DPP achieves the highest coverage because it optimizes diversity exclusively without quality weighting; DivFlow trades some coverage for substantially better response quality.
All synthetic comparisons are significant after Bonferroni correction ($m=3$).

\subsection{Composition Theorem Verification}

\begin{table}[t]
\centering
\caption{Algebraic proof verification results across 100 perturbation tests.}
\label{tab:composition}
\begin{tabular}{lcc}
\toprule
\textbf{Property} & \textbf{Max Error} & \textbf{Status} \\
\midrule
Welfare decomposition $\Delta W = \lambda \Delta q_i$ & $2.57 \times 10^{-16}$ & PASS \\
Plan reconstruction & $5.55 \times 10^{-17}$ & PASS \\
Payment independence & $0.0$ & PASS \\
\bottomrule
\end{tabular}
\end{table}

All six verification conditions pass at or below machine precision (Table~\ref{tab:composition}), confirming the algebraic proof.

\subsection{IC Violation Analysis}

\begin{table}[t]
\centering
\caption{IC violation analysis with VCG condition identification. 1{,}200 total tests (200 random + 1{,}000 adversarial) on 20-agent instances. Clopper--Pearson 95\% CIs.}
\label{tab:ic}
\begin{tabular}{lccc}
\toprule
\textbf{VCG Condition} & \textbf{Max Error} & \textbf{Status} \\
\midrule
C1: Quasi-linearity & $3.33 \times 10^{-16}$ & SATISFIED \\
C2: Exact maximization & greedy approx.\ & VIOLATED \\
C3: Payment correctness & $0.0$ & SATISFIED \\
\midrule
\multicolumn{3}{l}{\textbf{Violation taxonomy} (318 violations total)} \\
\midrule
Type A (selection boundary) & 318/318 & 100\% \\
Type B (payment distortion) & 0/318 & 0\% \\
Type C (submodularity slack) & 0/318 & 0\% \\
\midrule
Overall violation rate & \multicolumn{2}{c}{26.5\%} \\
Max utility gain & \multicolumn{2}{c}{0.56} \\
$\varepsilon_{\mathrm{IC}}$ bound (corrected) & \multicolumn{2}{c}{7.92 (validated)} \\
\bottomrule
\end{tabular}
\end{table}

\subsection{Scoring Rule Properness}

\begin{table}[t]
\centering
\caption{Scoring rule properness (500 tests per rule, 3--8 outcomes). Clopper--Pearson 95\% CIs.}
\label{tab:properness}
\begin{tabular}{lccc}
\toprule
\textbf{Rule} & \textbf{Violations} & \textbf{Rate} & \textbf{95\% CI} \\
\midrule
Logarithmic & 0/500 & 0.0\% & $[0.0\%,\; 0.7\%]$ \\
Brier & 0/500 & 0.0\% & $[0.0\%,\; 0.7\%]$ \\
Spherical & 0/500 & 0.0\% & $[0.0\%,\; 0.7\%]$ \\
CRPS & 0/500 & 0.0\% & $[0.0\%,\; 0.7\%]$ \\
Energy-augmented & 0/500 & 0.0\% & $[0.0\%,\; 0.7\%]$ \\
\bottomrule
\end{tabular}
\end{table}

All five scoring rules are verified proper with 0 violations across 500 tests each (Table~\ref{tab:properness}).

\subsection{Coverage Certificates}

Coverage certificates use Rogers' covering constant $C_{\mathrm{cov}} = 3$:
\[
N(\varepsilon,\, \mathcal{M}) \leq \left(\frac{3D}{\varepsilon}\right)^{d_{\mathrm{eff}}}.
\]
Using effective dimension $d_{\mathrm{eff}}$ (estimated via PCA at 95\% variance) rather than ambient dimension makes the bound non-trivial.

\subsection{Sensitivity Analysis}

We conduct a systematic sensitivity analysis across key hyperparameters (Table~\ref{tab:sensitivity}).

\begin{table}[t]
\centering
\caption{Hyperparameter sensitivity analysis. Violation rate (with 95\% CI) and topic coverage as a function of Sinkhorn regularization $\varepsilon$, quality weight $\lambda$, and selection size $k$.}
\label{tab:sensitivity}
\small
\begin{tabular}{lccc}
\toprule
\textbf{Parameter} & \textbf{Value} & \textbf{Viol.\ Rate [95\% CI]} & \textbf{Topic Cov.} \\
\midrule
\multicolumn{4}{l}{\textit{Sinkhorn $\varepsilon$ (quality weight $\lambda=0.3$, $k=3$)}} \\
$\varepsilon$ & 0.01 & 2.0\% $[0.2\%,\, 7.0\%]$ & 75.0\% \\
$\varepsilon$ & 0.10 & 14.0\% $[7.9\%,\, 22.4\%]$ & 75.0\% \\
$\varepsilon$ & 1.00 & 14.0\% $[7.9\%,\, 22.4\%]$ & 75.0\% \\
\midrule
\multicolumn{4}{l}{\textit{Quality weight $\lambda$ ($\varepsilon=0.1$, $k=3$)}} \\
$\lambda$ & 0.0 & 4.0\% $[1.1\%,\, 9.9\%]$ & 75.0\% \\
$\lambda$ & 0.3 & 14.0\% $[7.9\%,\, 22.4\%]$ & 75.0\% \\
$\lambda$ & 1.0 & 0.0\% $[0.0\%,\, 7.1\%]$ & 50.0\% \\
\midrule
\multicolumn{4}{l}{\textit{Selection size $k$ ($\varepsilon=0.1$, $\lambda=0.3$)}} \\
$k$ & 3 & 14.0\% & 75.0\% \\
$k$ & 5 & 20.0\% & 100.0\% \\
$k$ & 10 & 26.0\% & 100.0\% \\
\bottomrule
\end{tabular}
\end{table}

Key findings: (1)~Lower Sinkhorn $\varepsilon$ reduces violation rates (2.0\% at $\varepsilon=0.01$ vs.\ 14.0\% at $\varepsilon=0.10$) because tighter regularization improves submodularity.
(2)~$\lambda=0$ (pure diversity) has lower violations but at $\lambda=1$ (pure quality), violations vanish because the welfare function becomes exactly linear.
(3)~Larger $k$ increases violations because more agents are near the selection boundary, but also achieves 100\% topic coverage.

%% ===================================================================
\section{Related Work}
\label{sec:related}

\paragraph{Diverse text generation.}
\citet{li2023diverse} propose diverse beam search; \citet{wang2023selfconsistency} use self-consistency with majority voting.
These modify generation; DivFlow operates post-hoc.

\paragraph{DPP for text.}
\citet{kulesza2012determinantal} survey DPPs; \citet{cho2019multi} apply DPPs to summarization.

\paragraph{Optimal transport in ML.}
Sinkhorn divergence~\citep{cuturi2013sinkhorn, feydy2019interpolating} has been used for generative modeling.
We leverage Sinkhorn dual potentials for greedy subset selection with mechanism design integration.

\paragraph{Mechanism design.}
VCG mechanisms~\citep{vickrey1961counterspeculation, groves1973incentives} are classical.
\citet{nemhauser1978analysis} established the $(1-1/e)$ approximation for greedy submodular maximization.
Our composition theorem provides the first algebraic proof of Sinkhorn-VCG compatibility.

\paragraph{SMT verification.}
Z3~\citep{demoura2008z3} has been applied to verification of auction mechanisms~\citep{brandt2017rolling}.
We provide the first SMT encoding of IC conditions for optimal transport-based mechanisms.

%% ===================================================================
\section{Conclusion}
\label{sec:conclusion}

We presented DivFlow, combining Sinkhorn divergence-guided selection with VCG mechanism design.
Our main technical contributions are:
(1)~an \emph{algebraic} proof that Sinkhorn-based welfare is exactly quasi-linear (verified through the full welfare pipeline, not tautologically), with max decomposition error $2.57 \times 10^{-16}$;
(2)~a corrected $\varepsilon$-IC bound for approximately submodular functions: $\varepsilon_{\mathrm{IC}} \leq (1-\alpha_g)W(S^*) + k\delta_{\mathrm{sub}}$, which properly contains empirical violations (prior bound $W(S^*)/e$ was too tight);
(3)~Z3 SMT verification with grid-certified vs.\ soundness-certified distinction, and Lipschitz-based gap analysis;
(4)~evaluation on real LLM outputs (20 prompts, \texttt{gpt-4.1-nano} + \texttt{text-embedding-3-small}), showing DivFlow achieves the Pareto-optimal quality--diversity tradeoff;
(5)~a genuine DPP baseline (pure diversity kernel, not quality-weighted) that properly differs from TopQuality.
All results are generated from a single reproducible pipeline run with consistent numerical values throughout.

%% ===================================================================
\bibliographystyle{plainnat}
\begin{thebibliography}{20}

\bibitem[Brandt et~al.(2017)]{brandt2017rolling}
F.~Brandt, V.~Conitzer, U.~Endriss, J.~Lang, and A.~D. Procaccia, editors.
\newblock \emph{Handbook of Computational Social Choice}.
\newblock Cambridge University Press, 2017.

\bibitem[Carbonell and Goldstein(1998)]{carbonell1998mmr}
J.~Carbonell and J.~Goldstein.
\newblock The use of MMR, diversity-based reranking for reordering documents and producing summaries.
\newblock In \emph{SIGIR}, 1998.

\bibitem[Cho et~al.(2019)]{cho2019multi}
S.~Cho, L.~Lebanoff, H.~Foroosh, and F.~Liu.
\newblock Multi-document summarization with determinantal point processes.
\newblock In \emph{ACL Workshop}, 2019.

\bibitem[Clarke(1971)]{clarke1971multipart}
E.~H. Clarke.
\newblock Multipart pricing of public goods.
\newblock \emph{Public Choice}, 11:17--33, 1971.

\bibitem[Cuturi(2013)]{cuturi2013sinkhorn}
M.~Cuturi.
\newblock Sinkhorn distances: Lightspeed computation of optimal transport.
\newblock In \emph{NeurIPS}, 2013.

\bibitem[de~Moura and Bj{\o}rner(2008)]{demoura2008z3}
L.~de~Moura and N.~Bj{\o}rner.
\newblock Z3: An efficient SMT solver.
\newblock In \emph{TACAS}, 2008.

\bibitem[Feydy et~al.(2019)]{feydy2019interpolating}
J.~Feydy, T.~S\'{e}journ\'{e}, F.-X.~Vialard, S.-i.~Amari, and A.~Trouve.
\newblock Interpolating between optimal transport and MMD using Sinkhorn divergences.
\newblock In \emph{AISTATS}, 2019.

\bibitem[Groves(1973)]{groves1973incentives}
T.~Groves.
\newblock Incentives in teams.
\newblock \emph{Econometrica}, 41(4):617--631, 1973.

\bibitem[Kulesza and Taskar(2012)]{kulesza2012determinantal}
A.~Kulesza and B.~Taskar.
\newblock Determinantal point processes for machine learning.
\newblock \emph{Found.\ Trends Mach.\ Learn.}, 5(2--3):123--286, 2012.

\bibitem[Li et~al.(2023)]{li2023diverse}
Y.~Li, Z.~Lin, S.~Zhang, et~al.
\newblock Making language models better reasoners with step-aware verifier.
\newblock In \emph{ACL}, 2023.

\bibitem[Nemhauser et~al.(1978)]{nemhauser1978analysis}
G.~L. Nemhauser, L.~A. Wolsey, and M.~L. Fisher.
\newblock An analysis of approximations for maximizing submodular set functions.
\newblock \emph{Math.\ Programming}, 14(1):265--294, 1978.

\bibitem[Peyr\'{e} and Cuturi(2019)]{peyre2019computational}
G.~Peyr\'{e} and M.~Cuturi.
\newblock Computational optimal transport.
\newblock \emph{Found.\ Trends Mach.\ Learn.}, 11(5-6):355--607, 2019.

\bibitem[Vickrey(1961)]{vickrey1961counterspeculation}
W.~Vickrey.
\newblock Counterspeculation, auctions, and competitive sealed tenders.
\newblock \emph{J.\ Finance}, 16(1):8--37, 1961.

\bibitem[Wang et~al.(2023)]{wang2023selfconsistency}
X.~Wang, J.~Wei, D.~Schuurmans, et~al.
\newblock Self-consistency improves chain of thought reasoning in language models.
\newblock In \emph{ICLR}, 2023.

\end{thebibliography}

%% ===================================================================
%% APPENDIX
%% ===================================================================
\appendix
\newpage

\section{Full Algorithm Details}
\label{app:algorithms}

\subsection{Sinkhorn-Knopp Algorithm}

\begin{algorithm}[H]
\caption{Sinkhorn-Knopp for Entropic OT}
\label{alg:sinkhorn}
\begin{algorithmic}[1]
\REQUIRE Source marginal $a \in \Delta^{n-1}$, target marginal $b \in \Delta^{m-1}$, cost matrix $M \in \mathbb{R}^{n \times m}$, regularization $\varepsilon > 0$, iterations $T$
\ENSURE Dual potentials $(f, g)$
\STATE $K \leftarrow \exp(-M/\varepsilon)$
\STATE $u \leftarrow \mathbf{1}_n$, $v \leftarrow \mathbf{1}_m$
\FOR{$t = 1, \ldots, T$}
  \STATE $u \leftarrow a \oslash (K v)$
  \STATE $v \leftarrow b \oslash (K^\top u)$
  \IF{$\|u - u_{\mathrm{prev}}\|_\infty < \mathrm{tol}$}
    \STATE \textbf{break}
  \ENDIF
\ENDFOR
\STATE $f \leftarrow \varepsilon \log u$, $g \leftarrow \varepsilon \log v$
\RETURN $(f, g)$
\end{algorithmic}
\end{algorithm}

\subsection{VCG Payment Computation}

\begin{algorithm}[H]
\caption{VCG Payments for Diverse Selection}
\label{alg:vcg}
\begin{algorithmic}[1]
\REQUIRE Embeddings $\{x_i\}$, qualities $\{q_i\}$, selected set $S^*$, welfare $W$
\ENSURE Payments $\{p_i\}_{i \in S^*}$
\FOR{each $i \in S^*$}
  \STATE $W_{-i}^* \leftarrow W(S^* \setminus \{i\})$ \COMMENT{welfare of others in $S^*$}
  \STATE $S_{-i} \leftarrow \argmax_{|S|=k,\, i \notin S} W(S)$ \COMMENT{optimal without $i$}
  \STATE $p_i \leftarrow \max(W(S_{-i}) - W_{-i}^*,\; 0)$ \COMMENT{externality}
\ENDFOR
\RETURN $\{p_i\}$
\end{algorithmic}
\end{algorithm}

\section{Detailed Algebraic Proof}
\label{app:proof}

We provide the full algebraic derivation of Theorem~\ref{thm:composition}.

\subsection{Sinkhorn Dual Representation}

The entropic optimal transport dual is:
\[
\mathrm{OT}_\varepsilon(\mu, \nu) = \max_{f, g} \sum_i a_i f_i + \sum_j b_j g_j - \varepsilon \sum_{i,j} a_i b_j \left[\exp\!\left(\frac{f_i + g_j - C_{ij}}{\varepsilon}\right) - 1\right].
\]
At optimality, the dual potentials satisfy the Sinkhorn fixed-point equations:
\begin{align}
f_i &= \varepsilon \log a_i - \varepsilon \log \sum_j \exp\!\left(\frac{g_j - C_{ij}}{\varepsilon}\right), \\
g_j &= \varepsilon \log b_j - \varepsilon \log \sum_i \exp\!\left(\frac{f_i - C_{ij}}{\varepsilon}\right).
\end{align}

\subsection{Independence Argument}

\begin{enumerate}
\item The marginals $a_i = 1/|S|$ depend on $|S|$ but not on $q_i$.
\item The cost matrix $C_{ij} = \|x_i - y_j\|^2$ depends on embeddings, not qualities.
\item The Gibbs kernel $K_{ij} = \exp(-C_{ij}/\varepsilon)$ depends on embeddings, not qualities.
\item The Sinkhorn iterates $(u, v)$ are functions of $K$, $a$, and $b$ only.
\item The dual potentials $f = \varepsilon \log u$ and $g = \varepsilon \log v$ are thus quality-independent.
\item The optimal plan $\pi^*_{ij} = \exp((f_i + g_j - C_{ij})/\varepsilon)$ is quality-independent.
\item The Sinkhorn divergence $\mathcal{S}_\varepsilon(\mu, \nu)$ is quality-independent.
\end{enumerate}

Therefore $\mathrm{div}(S) = -\mathcal{S}_\varepsilon(\mu_S, \nu)$ is independent of all $\{q_i\}$, and the decomposition $W(S) = h_i + \lambda q_i \mathbf{1}[i \in S]$ is exact.

\subsection{Numerical Verification}

We verify each step:
\begin{itemize}
\item Steps 1--6: The algebraic proof verifies quality independence through the full welfare pipeline: perturbing $q_i$ and checking that $\Delta W = \lambda \Delta q_i$ exactly. Max decomposition error: $2.57 \times 10^{-16}$.
\item Step 7: Max $|W_{\text{actual}} - W_{\text{predicted}}| = 2.57 \times 10^{-16}$.
\item VCG payment: Max $|p_i(q_i) - p_i(q_i')| = 0$ across all perturbations.
\end{itemize}

\section{IC Violation Root-Cause Analysis: Full Details}
\label{app:ic}

\subsection{VCG Condition Failure Mechanism}

When $S^*$ is the true welfare maximizer and $S_g$ the greedy solution:
\begin{enumerate}
\item If $S_g = S^*$: VCG is DSIC. No violations possible.
\item If $S_g \neq S^*$: The VCG payment $p_i = W_{-i}(S_{g,-i}) - W_{-i}(S_g)$ uses $S_g$ instead of $S^*$. An agent can potentially exploit the gap by misreporting quality to change the greedy selection order.
\end{enumerate}

\subsection{Violation Bound Derivation}

For \emph{exactly} submodular functions, the Nemhauser--Wolsey--Fisher theorem gives:
\[
W(S_g) \geq (1 - 1/e) \cdot W(S^*).
\]
However, Sinkhorn divergence is only $\delta_{\mathrm{sub}}$-approximately submodular.
The corrected maximum IC violation bound is:
\[
\varepsilon_{\mathrm{IC}} \leq (1 - \alpha_g) \cdot W(S^*) + k \cdot \delta_{\mathrm{sub}},
\]
where $\alpha_g$ is the empirical greedy approximation ratio.

For our experimental instance with $\delta_{\mathrm{sub}} = 1.49$ and $k = 5$, the corrected bound is $7.92$, which properly contains the observed maximum utility gain of $0.56$.

\subsection{Type A Violation Mechanism}

Type A violations occur when:
\begin{enumerate}
\item Agent $i$ is not in $S_g$ (not selected under truthful reporting).
\item By inflating $q_i$, agent $i$ enters the greedy selection.
\item Agent $i$'s true quality minus the VCG payment yields positive utility.
\end{enumerate}

This is more likely when agent $i$'s marginal welfare gain is close to the selection threshold.

\section{Z3 Encoding Details}
\label{app:z3}

\subsection{SMT Formulation}

For each agent $i$ and pair of allocations $(S_{\text{truth}}, S_{\text{dev}})$:

\begin{verbatim}
; Symbolic variables
(declare-const q_true Real)
(declare-const q_fake Real)
(assert (>= q_true 0.0))
(assert (<= q_true 1.0))
(assert (>= q_fake 0.0))
(assert (<= q_fake 1.0))
(assert (not (= q_fake q_true)))

; Welfare under truth (div_S is pre-computed constant)
; W(S_truth, q_true) = div_S + lambda * (sum_others + q_true)
(define-fun w_truth () Real
  (+ div_truth (* lambda (+ sum_others_truth q_true))))

; Optimality constraints (S_truth maximizes W under q_true)
; ... for all other subsets S' ...

; IC violation: u_dev > u_truth + epsilon
(assert (> (- q_true p_dev) (+ (- q_true p_truth) epsilon)))
\end{verbatim}

Z3 then searches for satisfying assignments, producing concrete counterexamples.

\section{Statistical Methodology}
\label{app:stats}

\subsection{Enhanced Bootstrap}

We use $n_{\text{seeds}} = 20$ independent bootstrap runs, each with $n_{\text{boot}} = 2000$ resamples, for a total of 40{,}000 bootstrap samples.
CI stability is assessed by the standard deviation of CI endpoints across seeds.

\subsection{Multiple Testing Correction}

With $m$ hypothesis tests at level $\alpha$:
\begin{itemize}
\item \textbf{Bonferroni}: Reject $H_i$ if $p_i < \alpha/m$. Controls family-wise error rate.
\item \textbf{Benjamini--Hochberg}: Sort $p_{(1)} \leq \cdots \leq p_{(m)}$. Find largest $k$ s.t.\ $p_{(k)} \leq k\alpha/m$. Reject $H_{(1)}, \ldots, H_{(k)}$. Controls false discovery rate.
\end{itemize}

In our experiments with 4 comparisons at $\alpha = 0.05$:
\begin{itemize}
\item Bonferroni threshold: $0.05/4 = 0.0125$.
\item BH at FDR 0.05: adaptive thresholds $0.0125, 0.025, 0.0375, 0.05$.
\end{itemize}

\subsection{Power Analysis}

For detecting a difference in proportions $p_0$ vs.\ $p_1$:
\[
n \geq \frac{(z_{\alpha/2} + z_\beta)^2 \cdot (p_0(1-p_0) + p_1(1-p_1))}{(p_0 - p_1)^2}.
\]

\section{Coverage Certificates: Full Derivation}
\label{app:coverage}

\subsection{Rogers' Covering Constant}

The Rogers bound states: a maximal $\varepsilon$-packing of $B_d(R)$ has size $\leq (2R/\varepsilon + 1)^d$, and any maximal $\varepsilon$-packing is a $2\varepsilon$-covering.
This gives $N(\varepsilon) \leq (3R/\varepsilon)^d$ with $C_{\text{cov}} = 3$.

\subsection{Effective Dimension Estimation}

We estimate $d_{\text{eff}}$ as the number of principal components explaining 95\% of total variance.
For LLM embeddings in ambient $d = 1536$, typical $d_{\text{eff}} \approx 4$--$8$, making metric entropy bounds non-trivial.

\section{Scoring Rule Properness Proofs}
\label{app:properness}

\begin{proposition}[Energy-augmented properness]
\label{prop:energy-proper}
If $S_{\text{base}}$ is strictly proper, then $S_E(p, y) = S_{\text{base}}(p, y) + \lambda \cdot E(y, Y_{\text{hist}})$ is strictly proper.
\end{proposition}

\begin{proof}
$\mathbb{E}_q[S_E(p, Y)] = \mathbb{E}_q[S_{\text{base}}(p, Y)] + \lambda \cdot \mathbb{E}_q[E(Y, Y_{\text{hist}})]$.
The second term is constant w.r.t.\ $p$, so $\argmax_p \mathbb{E}_q[S_E(p, Y)] = \argmax_p \mathbb{E}_q[S_{\text{base}}(p, Y)] = q$.
\end{proof}

\section{Experimental Details and Reproducibility}
\label{app:reproduce}

All experiments use NumPy random seeds for reproducibility.
Code is available at \url{https://github.com/anonymous/divflow}.

\paragraph{Compute.} All experiments run on a single CPU in under 10 minutes.
No GPU is required.

\paragraph{Software.} Python 3.9+, NumPy 1.22+, SciPy 1.9+, Z3 4.12+.

\end{document}
